% ID: SEL-D0027-OB-E04
% ESTADO: cotejada. Segunda pasada visual completa y borrador compilado revisados.
% Fuente: Selectividad/matematicas_aplicadas/suplemente_septiembre_EXAMEN 6-ANDALUCÍA.pdf
% SHA-256: 601fd1292e7d6b21ceecde2856a283e0e929d02d4e314d97fd7062bda846650a
% Página: 2; opción B; ejercicio original 4.

El precio de un determinado producto se distribuye según una ley Normal de desviación típica 5 \texteuro{} y media desconocida. Se toman 10 comercios al azar y se observa en ellos el precio de este producto, resultando los siguientes valores en euros:
\[
96\quad108\quad97\quad112\quad99\quad106\quad105\quad100\quad98\quad99
\]
\begin{enumerate}[label=\alph*)]
\item \textbf{(0.5 puntos)} \textquestiondown{}Cuál es la distribución del precio medio del producto en las muestras de tamaño 10?
\item \textbf{(1 punto)} Determine un intervalo de confianza, al 97\%, para la media poblacional.
\item \textbf{(1 punto)} Con el mismo nivel de confianza, \textquestiondown{}cuál debe ser el tamaño mínimo de la muestra de esa población para que el error cometido sea menor que 2?
\end{enumerate}
